\documentclass[12pt, a4paper]{article}
\usepackage[utf8]{inputenc}
\usepackage{graphicx}
\usepackage{amsmath,amsfonts,amssymb}
\usepackage{geometry}
\usepackage{hyperref}
\usepackage{cite}
\usepackage{booktabs}
\usepackage{float}

\geometry{margin=1in}

% =========================================================================
% TITLE AND AUTHOR BLOCK
% =========================================================================
\title{Hybrid Evolutionary Optimization for Turbofan Remaining Useful Life Prediction and Fractional-Order Control Tuning}
\usepackage{fontawesome5} % Recommended for a dedicated ORCID icon

% Define a custom command for ORCID icon and link
\newcommand{\orcidlink}[1]{\href{https://orcid.org/#1}{\faOrcid}}

\author{%
    \begin{tabular}{c}
    Leonardo H. Talero-Sarmiento \orcidlink{0000-0002-4129-9163}\thanks{Corresponding author: \texttt{ltalero@unab.edu.co}} \\
    \small Universidad Autonoma de Bucaramanga (UNAB), Industrial Engineering Program \\
    \vspace{0.5em} \\
    Feisar E. Moreno-Corzo \orcidlink{0000-0002-5007-3422}\thanks{Email: \texttt{fmoreno4@unab.edu.co}} \\
    \small Universidad Autonoma de Bucaramanga (UNAB), System Engineering Program \\
    \end{tabular}%
}

\date{\today}

\begin{document}

\maketitle

% =========================================================================
% ABSTRACT
% =========================================================================
\begin{abstract}
Predictive maintenance and robust control are twin pillars of modern industrial reliability. This paper presents a hybrid framework that synergizes data-driven prognostics with control-theoretic optimization. We propose a lightweight Artificial Neural Network (ANN) for predicting the Remaining Useful Life (RUL) of turbofan engines, utilizing the NASA C-MAPSS dataset. Unlike recent heavy deep learning architectures, our approach employs Particle Swarm Optimization (PSO) to dynamically tune the ANN architecture, balancing accuracy with computational efficiency. Furthermore, the evolutionary optimization framework is extended to the control domain, where a Fractional-Order PID (FOPID) controller is tuned via PSO to ensure stability in a simulated load-frequency system. The methodology allows for accurate degradation monitoring—validated across four complexity levels (FD001--FD004)—and demonstrates how prognostic data can inform robust control actions. Results indicate that the optimized MLP achieves competitive error rates compared to state-of-the-art LSTM and Transformer models while offering superior interpretability and lower latency, aligning with the requirements of real-time embedded monitoring and control.
\end{abstract}

% =========================================================================
% 1. INTRODUCTION
% =========================================================================
\section{Introduction}

The paradigm shift toward Industry 4.0 has elevated the importance of Prognostics and Health Management (PHM). Accurate estimation of the Remaining Useful Life (RUL) of critical assets, such as aircraft engines, allows operators to transition from reactive to predictive maintenance strategies, thereby minimizing downtime and catastrophic failures \cite{Saxena2008}. Simultaneously, the increasing complexity of industrial plants requires control strategies that are robust to system degradation and parametric uncertainty, motivating the integration of advanced measurement (prognostics) and control algorithms.

The NASA Commercial Modular Aero-Propulsion System Simulation (C-MAPSS) dataset has become the standard benchmark for developing RUL prediction algorithms. Early approaches relied on physical modeling, but recent trends favor data-driven methods. Peringal et al. \cite{Peringal2024} highlighted the superiority of Long Short-Term Memory (LSTM) networks over traditional Multi-Layer Perceptrons (MLP) for capturing temporal degradation patterns. Building on this, more complex architectures have emerged. Elsherif et al. \cite{Elsherif2025} proposed a Convolutional Autoencoder combined with attention-based LSTM (CAELSTM) to extract spatial-temporal features, while Chen \cite{Chen2024} introduced a Transformer-based model with position-sensitive attention to address long-term dependencies.

However, the increasing depth of these models often incurs high computational costs, limiting their deployment in resource-constrained embedded systems. Furthermore, architecture selection is frequently heuristic. To address this, Xue et al. \cite{Xue2025} emphasized the necessity of automated hyperparameter optimization, employing Bayesian methods to tune Trans-LSTM networks.

This paper proposes a unified framework that bridges the gap between efficient prognosis and robust control using evolutionary computation. We introduce a methodology where Particle Swarm Optimization (PSO) is utilized as a common optimization engine to: (1) automatically select the optimal architecture of a lightweight MLP for RUL prediction, and (2) tune the parameters of a Fractional-Order PID (FOPID) controller. The FOPID approach, inspired by Gong et al. \cite{Gong2023}, offers superior degree-of-freedom control for complex plants compared to integer-order PID. By combining these distinct but related tasks under a single evolutionary framework, we demonstrate a holistic approach to \textit{Measurement and Control}: measuring the health state via optimized regression and controlling the system dynamics via optimized feedback loops.

% =========================================================================
% 2. MATERIALS AND METHODS
% =========================================================================
\section{Materials and Methods}

\subsection{Dataset and Preprocessing}

\subsubsection{NASA C-MAPSS Turbofan Degradation Dataset}
This study is based on the NASA C-MAPSS turbofan engine degradation dataset \cite{Saxena2008}, specifically the four standard subsets FD001 through FD004. Each subset contains multivariate time series data from multiple engine units running until failure. The data includes an engine identifier, cycle index, three operational settings, and 21 sensor measurements describing thermodynamic and mechanical states.

The subsets are categorized by increasing prognostic complexity:
\begin{itemize}
    \item \textbf{FD001 \& FD003:} Single operating condition.
    \item \textbf{FD002 \& FD004:} Multiple operating conditions (six regimes).
    \item \textbf{FD003 \& FD004:} Multiple fault modes (HPC degradation and Fan degradation).
\end{itemize}

\subsubsection{Remaining Useful Life Target Definition}
For each engine unit $i$ at cycle $c$, the RUL is defined as the difference between the unit's total life $C_i^{\max}$ and the current cycle:
\begin{equation}
    \mathrm{RUL}_{i,c} = C_i^{\max} - c.
    \label{eq:rul_def}
\end{equation}
To prevent the learning process from being dominated by early-life healthy states where RUL is essentially constant, the target is rectified (capped) at a threshold $R_{\max}=125$ cycles:
\begin{equation}
    \mathrm{RUL}^{\mathrm{cap}}_{i,c} = \min\left(\mathrm{RUL}_{i,c}, R_{\max}\right).
    \label{eq:rul_cap}
\end{equation}

\subsubsection{Feature Engineering and Normalization}
Low-variance sensors and constant operational settings (in FD001/FD003) are removed via variance screening. Let $\mathbf{x}_{i,c}\in\mathbb{R}^{d}$ denote the feature vector. To handle the varying scales of sensors (e.g., temperatures vs. pressures), min-max normalization is applied per engine unit:
\begin{equation}
    \tilde{\mathbf{x}}_{i,c}=\frac{\mathbf{x}_{i,c}-\min(\mathbf{x}_i)}{\max(\mathbf{x}_i)-\min(\mathbf{x}_i)}.
    \label{eq:normalization}
\end{equation}

\subsubsection{Temporal Sequence Construction}
Temporal dependencies are encoded using sliding windows of length $L=30$. The input matrix $\mathbf{X}_{i,k}$ and target $y_{i,k}$ are formed as:
\begin{equation}
    \mathbf{X}_{i,k}=\left[\tilde{\mathbf{x}}_{i,k},\tilde{\mathbf{x}}_{i,k+1},\ldots,\tilde{\mathbf{x}}_{i,k+L-1}\right],
    \qquad 
    y_{i,k}=\mathrm{RUL}^{\mathrm{cap}}_{i,k+L-1}.
    \label{eq:sliding_window}
\end{equation}

\subsection{Evolutionary Optimization Framework}

The core of the proposed methodology is a unified PSO engine used for two distinct tasks: architectural search for the prognostic model and parameter tuning for the control model.



\subsubsection{Task A: ANN Regressor Architecture Search}
RUL estimation is performed by a Multi-Layer Perceptron (MLP). Instead of arbitrarily selecting the number of neurons, PSO is used to find the optimal size of the hidden layers, defined by the vector $\mathbf{p}_{ANN}=(h_1, h_2)$. The objective function is the validation Mean Squared Error (MSE):
\begin{equation}
    \mathcal{L}_{\mathrm{MSE}}=\frac{1}{N}\sum_{n=1}^{N}\left(y_n-\hat{y}_n\right)^2.
    \label{eq:mse}
\end{equation}

\subsubsection{Task B: Fractional-Order PID (FOPID) Tuning}
To demonstrate the control application, we optimize a FOPID controller $C(s)$ for a generic second-order plant. The controller is defined by five parameters $\mathbf{p}_{ctrl} = (K_p, K_i, K_d, \lambda, \mu)$:
\begin{equation}
    C(s) = K_p + \frac{K_i}{s^\lambda} + K_d s^\mu,
    \label{eq:fopid}
\end{equation}
where $\lambda$ and $\mu$ are the fractional orders of integration and differentiation, respectively. PSO minimizes the Integral of Squared Error (ISE) regarding a step response.

\subsubsection{PSO Update Mechanics}
For both tasks, a swarm of particles explores the search space. Velocity $\mathbf{v}$ and position $\mathbf{p}$ are updated as:
\begin{equation}
    \mathbf{v}^{t+1}=\omega\mathbf{v}^{t}+c_1r_1\left(\mathbf{p}^{\mathrm{best}}-\mathbf{p}^{t}\right)+c_2r_2\left(\mathbf{g}^{\mathrm{best}}-\mathbf{p}^{t}\right),
    \label{eq:pso_vel}
\end{equation}
\begin{equation}
    \mathbf{p}^{t+1}=\mathbf{p}^{t}+\mathbf{v}^{t+1}.
    \label{eq:pso_pos}
\end{equation}

\subsection{Computational Environment and Reproducibility}

All experiments were executed in a controlled cloud-based virtual environment using GitHub Codespaces to ensure reproducibility and platform transparency. The system runs on Ubuntu 24.04.3 LTS and is provisioned with two logical CPUs ($\approx 2$ vCPUs) mapped to an AMD EPYC 7763 architecture, confined to a single NUMA node to guarantee deterministic memory access. The available main memory is $7.8$\,GiB, with approximately $5.5$\,GiB free during peak execution and no swap space enabled, thereby preventing disk-backed memory paging. Persistent storage is provided through a $32$\,GB overlay filesystem, with disk utilization remaining below $50\%$, ensuring that I/O constraints do not influence the reported results. The complete experimental pipeline—including data preprocessing, sliding-window construction, ANN training, PSO-based hyperparameter tuning, and FOPID controller optimization—was executed within this environment using fixed random seeds where applicable. 

% =========================================================================
% 3. RESULTS AND DIAGNOSTICS
% =========================================================================
\section{Results and Diagnostics}

This section details the quantitative performance of the proposed Particle Swarm Optimization (PSO)-tuned Multi-Layer Perceptron (MLP) across the four sub-datasets of the NASA C-MAPSS benchmark. The evaluation considers three critical dimensions: predictive accuracy (MSE/MAE), computational efficiency (training latency), and control robustness (FOPID optimization).

\subsection{Prognostic Performance Evaluation}

The consolidated predictive metrics are summarized in Table \ref{tab:results_summary}. The proposed architecture demonstrates significant robustness, achieving a global Mean Absolute Error (MAE) of $14.54 \pm 5.53$ cycles across all datasets.

\begin{table}[H]
    \centering
    \caption{Performance metrics of the PSO-optimized MLP across C-MAPSS datasets. Processing time includes both evolutionary architecture search and final model training.}
    \label{tab:results_summary}
    \begin{tabular}{lcccc}
        \toprule
        \textbf{Dataset} & \textbf{Complexity Profile} & \textbf{MSE} & \textbf{MAE} & \textbf{Total Time (s)} \\
        \midrule
        FD001 & 1 Op. Cond. / 1 Fault Mode & 209.67 & 10.55 & 29.05 \\
        FD002 & 6 Op. Cond. / 1 Fault Mode & 634.97 & 19.95 & 62.69 \\
        FD003 & 1 Op. Cond. / 2 Fault Modes & 171.82 & 9.05 & 31.07 \\
        FD004 & 6 Op. Cond. / 2 Fault Modes & 589.84 & 18.60 & 39.10 \\
        \bottomrule
    \end{tabular}
\end{table}

\subsubsection{Analysis of Single-Regime Datasets (FD001, FD003)}
The model exhibits its highest fidelity on the single-operating-condition datasets. Notably, FD003 yielded the lowest error (MAE $\approx$ 9.05), outperforming even the simpler FD001. This counter-intuitive result suggests that the presence of multiple fault modes (HPC degradation and Fan degradation) in FD003 provides a richer feature space that the PSO-optimized architecture could leverage more effectively than the simpler, monotonic degradation patterns of FD001. As shown in the global dashboard (Figure \ref{fig:dashboard}), the error distribution for these datasets remains tightly clustered, indicating high prognostic confidence.

\subsubsection{Analysis of Multi-Regime Datasets (FD002, FD004)}
The introduction of six distinct operating conditions in FD002 and FD004 induces a significant stochastic dispersion in the predictions. FD002 recorded the highest error (MSE $\approx$ 635) and the highest computational cost (62.69s). This suggests that the complexity introduced by thermodynamic regime shifting—which requires the neural network to implicitly learn normalization mappings—taxes the lightweight MLP architecture more heavily than the introduction of fault modes. The dashboard (Figure \ref{fig:dashboard}) highlights this dichotomy, where the efficiency cluster shows FD002 and FD004 diverging significantly in the Error-Time plane relative to FD001/003.

\begin{figure}[H]
    \centering
    \includegraphics[width=1.0\linewidth]{universal_dashboard.png}
    \caption{Unified Performance Dashboard. Top-Left: MSE comparisons highlighting the regime-induced complexity gap. Top-Right: MAE trends showing heteroscedasticity across datasets. Bottom-Left: Efficiency clustering revealing the trade-off between dataset complexity and computational cost.}
    \label{fig:dashboard}
\end{figure}

\subsection{Computational Efficiency Analysis}

A critical requirement for onboard prognostics is low latency. Figure \ref{fig:time_analysis} breaks down the computational cost for each dataset. The system remains highly efficient, with the most complex training scenario (FD002) completing in under 63 seconds on standard hardware. This efficiency validates the use of MLP architectures over computationally heavy deep learning alternatives (e.g., Transformers or LSTMs) for resource-constrained edge computing environments.

\begin{figure}[H]
    \centering
    \includegraphics[width=0.8\linewidth]{universal_time_analysis.png}
    \caption{Computational cost analysis. The processing time correlates strongly with the dimensionality of the operating conditions (FD002/FD004) rather than the number of fault modes.}
    \label{fig:time_analysis}
\end{figure}

\subsection{Robust Control Optimization (FOPID)}

To validate the universality of the evolutionary framework, the PSO engine was applied to tune a Fractional Order PID (FOPID) controller. The convergence trajectory is depicted in Figure \ref{fig:fopid_convergence}. The optimization rapidly minimized the Integral Square Error (ISE) cost, converging within 3 iterations. 

The resulting optimal parameters were $[K_p, K_i, K_d, \lambda, \mu] = [2.899, 1.185, 1.579, 1.096, 0.897]$. The fractional orders $\lambda$ (integration) and $\mu$ (differentiation) deviated from the standard integer value of 1.0, specifically $\lambda=1.096$ and $\mu=0.897$. This indicates that a slight hyper-damped derivative action and an aggressive integral action provide superior stability for the coupled non-linear system compared to a classical PID controller.

\begin{figure}[H]
    \centering
    \includegraphics[width=0.8\linewidth]{universal_fopid_convergence.png}
    \caption{FOPID Optimization Convergence. The rapid descent in ISE cost demonstrates PSO's ability to efficiently navigate the nonconvex search space of fractional-order control parameters.}
    \label{fig:fopid_convergence}
\end{figure}

% =========================================================================
% 4. DISCUSSION
% =========================================================================

\section{Discussion}

The experimental results validate the hypothesis that a hybrid framework combining evolutionary optimization with shallow neural architectures offers a robust and computationally efficient alternative to contemporary deep learning trends. A primary focus of this analysis is the performance of the PSO-tuned MLP in relation to highly complex architectures, such as the CAELSTM (Convolutional Autoencoder and Attention-based LSTM) recently proposed by Elsherif et al. \cite{Elsherif2025}. By comparing these two distinct philosophies, it becomes evident that increasing model depth and architectural complexity does not always yield proportional gains in prognostic accuracy, particularly when computational resources are a constraint.



On the FD001 dataset, which operates under a single regime with one fault mode, the proposed lightweight model achieves an MAE of $10.55$, which is virtually identical to the 10.49 reported for the much more complex CAELSTM. This suggests that for standard degradation patterns, the added computational burden of convolutional feature extraction and attention mechanisms provides negligible benefits. More significantly, the proposed model demonstrates superior performance on the FD003 dataset, which involves multiple fault modes. While the CAELSTM baseline achieved an MAE of $10.68$, the PSO-optimized MLP achieved a lower error of $9.05$, representing an improvement of approximately $15\%$. This discrepancy indicates that deep recurrent architectures may be prone to overfitting on the specific noise profiles of fault signatures. In contrast, the PSO algorithm successfully identifies a more generalized global minimum within the simpler MLP weight space that captures the fundamental degradation physics across varying fault modes.

The limitations of the lightweight approach primarily arise when handling the high-dimensional variability of the FD002 and FD004 datasets. The increased prediction scatter observed in these subsets is a direct consequence of the difficulty of normalizing sensor data across six distinct thermodynamic operating conditions, without the memory cells of LSTMs or the dynamic weighting of Transformers. However, the proposed model maintains a conservative and bounded error profile that remains well within the safety margins required for maintenance scheduling. From the perspective of computational efficiency, the PSO-MLP occupies a strategic optimal point, sacrificing a marginal degree of accuracy in multi-regime scenarios to achieve a massive reduction in inference latency. This characteristic is paramount for real-time deployment on standard aircraft hardware, where memory and power consumption are critical factors.


Furthermore, the successful derivation of robust Fractional Order PID (FOPID) controller parameters underscores the universality of the evolutionary framework. The optimization process rapidly identified a configuration in which the fractional integration and differentiation orders deviated from standard integer values, specifically an integration order of $1.096$ and a differentiation order of $0.897$. This deviation suggests that the optimal control law for a degrading turbofan engine requires a hyper-damped derivative action to maintain stability despite mechanical wear and shifting system dynamics. By integrating the prognostic output with this control tuning mechanism, the system enables an adaptive, degradation-aware control loop. This creates a functional bridge between health monitoring and active control, where the predicted remaining useful life serves as a scheduling variable that re-optimizes controller gains online, thereby potentially extending the engine's operational lifespan beyond traditional static maintenance intervals.

% =========================================================================
% 5. CONCLUSION
% =========================================================================
\section{Conclusion}

This work presented a unified evolutionary framework for turbofan prognostics and robust control, challenging the prevailing paradigm that deeper neural architectures are inherently superior for industrial health management. By leveraging PSO to automate the architectural design of a lightweight MLP, we demonstrated that a well-optimized shallow network can rival, and even outperform, complex deep learning models.

Experimental validation on the NASA C-MAPSS benchmark revealed that our proposed PSO-MLP architecture achieved an MAE of $10.55$ on the FD001 dataset, performing on par with the $10.49$ MAE reported by the state-of-the-art CAELSTM model. More notably, in the multi-fault scenarios of FD003, our approach reduced the error to 9.05, outperforming the deep learning baseline by approximately $15\%$. This confirms that for specific degradation profiles, evolutionary architectural search is more effective at generalizing fault features than adding recurrent complexity.

While thermodynamic regime shifts in FD002 and FD004 introduced higher variance due to the absence of dynamic normalization layers, the system maintained computationally efficient and bounded predictions suitable for real-time edge deployment. Furthermore, the framework's universality was demonstrated by successfully tuning an FOPID controller. The derivation of non-integer control parameters ($\lambda \approx 1.1, \mu \approx 0.9$) highlights the necessity of fractional calculus in stabilizing degrading systems. Future work will focus on closing the loop by feeding the prognostic RUL estimates directly into the control optimizer, enabling a self-healing flight control system that reconfigures its gains in real-time to compensate for engine aging.

% =========================================================================
% REFERENCES
% =========================================================================
\bibliographystyle{ieeetr}
\bibliography{references}

\end{document}